\begin{mistake}{2026-05-13}{05}

\begin{problemcontent}
设 \(y=y(x)\) 由参数方程
\[
\begin{cases} x = \dfrac{t}{1+t^3} \\[6pt] y = \dfrac{t^2}{1+t^3} \end{cases}
\]
确定，则曲线 \(y=y(x)\) 的斜渐近线方程为\underline{\hspace{3cm}}。
\end{problemcontent}

\begin{solutioncontent}
斜渐近线 \(y = kx + b\)，当 \(t \to -1\) 时 \(x \to \infty\)。

要求 \(\dfrac{t^2 - kt}{1+t^3}\) 在 \(t \to -1\) 有有限极限，分母为 0 则分子也须为 0：
\[
(-1)^2 - k(-1) = 1 + k = 0 \implies k = -1
\]

代入 \(k=-1\)，求 \(b\)：
\[
b = \lim_{t \to -1}\frac{t^2 + t}{1+t^3} = \lim_{t \to -1}\frac{t(t+1)}{(1+t)(1-t+t^2)} = \frac{-1}{3}
\]

斜渐近线方程为 \(y = -x - \dfrac{1}{3}\)。
\end{solutioncontent}

\mistakeknowledge{
\begin{itemize}
  \item \textbf{核心公式/定理}：斜渐近线 \(y=kx+b\)，其中 \(k=\lim_{x\to\infty}\frac{y}{x}\)，\(b=\lim_{x\to\infty}(y-kx)\)；参数方程下转化为对应参数极限。
  \item \textbf{易错点}：求 \(b\) 时 \(y - kx\) 中 \(k=-1\) 的符号代入易出错；因式分解 \(1+t^3=(1+t)(1-t+t^2)\) 需熟记。
  \item \textbf{关联知识}：参数方程中 \(x\to\infty\) 对应使分母为零的参数值；"迫使分子消零定系数"可快速确定 \(k\)。
\end{itemize}}

\end{mistake}
